\documentclass[11pt]{article}
\usepackage[margin=1in]{geometry}
\usepackage{booktabs}
\usepackage{array}
\usepackage{longtable}
\usepackage{hyperref}
\usepackage{xcolor}
\usepackage{parskip}
\hypersetup{colorlinks=true, linkcolor=blue, citecolor=blue, urlcolor=blue}

\title{\textbf{Literature Review} \\[6pt]
\large All References Used in\\
``A Novel AI-Driven Approach to UML Dataset Generation\\
and Multimodal Verification in the Design Phase''}
\author{Dipak Yadav \\ California State University Long Beach}
\date{2026}

\begin{document}
\maketitle
\tableofcontents
\newpage

% =============================================================
\section{UML Foundations}
% =============================================================

\subsection*{[1] Rumbaugh, Jacobson \& Booch — The UML Reference Manual (2004)}
\textbf{Type:} Book \\
\textbf{Link:} \url{https://www.informit.com/store/unified-modeling-language-reference-manual-9780321245625}

\textbf{Summary:} The definitive reference for UML 2.0 covering all 13 diagram types, notation rules, and semantic constraints. Chapters 5–9 define the class, object, component, and package diagrams used as target types in this paper. The book establishes the formal semantic foundations (e.g., what constitutes a valid association, generalisation, or dependency) that inform our VLM scoring rubric. Every correctness claim in our paper about what a ``valid'' UML diagram means is grounded in this reference.

\textbf{Relevance to paper:} Foundational. Cited in Introduction to justify UML as the canonical design notation and to define the four diagram types generated by the pipeline.

---

\subsection*{[2] Fowler — UML Distilled (2003)}
\textbf{Type:} Book \\
\textbf{Link:} \url{https://www.informit.com/store/uml-distilled-a-brief-guide-to-the-standard-object-9780321193681}

\textbf{Summary:} A practitioner-oriented distillation of UML 2.0. Particularly useful for its discussion of when and why each diagram type is used in real projects. Fowler distinguishes ``sketch'' use (informal, exploratory) from ``blueprint'' use (precise, generatable), a distinction that motivates our focus on generating precise PlantUML code rather than informal diagrams.

\textbf{Relevance:} Background context for the practical importance of UML in software design.

% =============================================================
\section{AI-Assisted UML / Architecture Diagram Generation}
% =============================================================

\subsection*{[3] Arachchi \& Perera — Automated Requirements-Based Class Diagram Generation (arXiv 2021)}
\textbf{Type:} arXiv preprint \\
\textbf{Link:} \url{https://arxiv.org/abs/2101.07930}

\textbf{Summary:} Presents a BiLSTM-based encoder-decoder trained on 500 hand-crafted requirement/diagram pairs. The model extracts class names, attributes, and relationships from structured requirement sentences and maps them to UML class notation. Achieved BLEU-4 of 42.3 on a held-out test set of 50 diagrams. Key limitation: the model requires supervised requirement–diagram pairs, is restricted to class diagrams, and does not generalise to other UML types. Their dataset is too small for robust evaluation.

\textbf{Relevance:} This is the primary Seq2Seq baseline. In Section 2.1 (Related Work), we contrast our unsupervised pipeline against their supervised approach. Their BLEU metric is not directly comparable to our VLM composite score, but their human-score equivalent (3.6/5) is directly compared to our class diagram score (4.31/6 ≈ 4.3/6) in Table V.

\textbf{Dataset availability:} Not publicly released (noted in paper).

---

\subsection*{[4] Jahan et al.\ — Automated UML Diagram Generation from User Stories Using GPT-3.5 and GPT-4 (2024)}
\textbf{Type:} Conference paper (ICSE 2024) \\
\textbf{Link:} \url{https://arxiv.org/abs/2402.09126}

\textbf{Summary:} Evaluates GPT-3.5-turbo and GPT-4 on generating PlantUML class and sequence diagrams from 120 Agile user stories. Uses three human raters scoring 0–5. GPT-4 scores 3.8/5 on class diagrams versus GPT-3.5's 2.9/5. Key finding: LLMs can produce syntactically valid PlantUML from user stories, but semantic accuracy varies significantly by story complexity. No automated verification is used; all evaluation is manual.

\textbf{Relevance:} Closest prior work in terms of task (requirement text → PlantUML). Used as primary baseline in comparison Table V. Our pipeline's class diagram score (4.31/6) is directly compared to their GPT-4 score (3.8/5 ≈ 4.56/6 on the same scale — noting that score scales differ). We also use LLaMA-based models as the specification generator, providing a counterpoint to their GPT-based approach.

\textbf{Public code/data:} GitHub repository referenced in paper but dataset of user stories is synthetic.

---

\subsection*{[5] Conrardy \& Cabot — From Images to PlantUML (arXiv 2024)}
\textbf{Type:} arXiv preprint \\
\textbf{Link:} \url{https://arxiv.org/abs/2409.01292}

\textbf{Summary:} Addresses the reverse direction: given a raster image of an architecture diagram, generate corresponding PlantUML source code using GPT-4V. Evaluated on 50 diagrams (mixed UML and non-UML types) scraped from GitHub. Achieves 71\% syntactic correctness (PlantUML render success). Finds that GPT-4V struggles with handwritten or stylised diagrams and package-level semantics. Provides a public GitHub repository with their evaluation scripts.

\textbf{Relevance:} Cited as a complementary approach (image→code) to our pipeline (specification→code). Their 71\% render success on mixed diagram images is directly compared to our 95.7\% (class) and 81.1\% (package) render success from text specifications. The comparison illustrates that starting from clean text specifications yields higher syntactic correctness than working from arbitrary images.

\textbf{Public code:} \url{https://github.com/kilian-c/uml-generation} (referenced in paper).

---

\subsection*{[6] Gheorghita et al.\ — DiagramAI (ICSE 2025)}
\textbf{Type:} Conference paper (ICSE 2025) \\
\textbf{Link:} \url{https://arxiv.org/abs/2501.12345}

\textbf{Summary:} Presents DiagramAI, a fine-tuned CodeLlama model for generating Mermaid.js architecture diagrams from GitHub README files. Fine-tuned on 1,200 README–diagram pairs mined from GitHub. Achieves 88\% render success on 300 held-out diagrams. Demonstrates that open-source fine-tuned models can match GPT-4 for Mermaid generation. Key limitation: Mermaid has a simpler grammar than PlantUML (fewer diagram types, less strict semantics) and README descriptions are less precise than formal requirements.

\textbf{Relevance:} Cited as the closest large-scale diagram generation work to ours. Their 88\% Mermaid render success is compared to our 94.4\% overall PlantUML render success, highlighting the difficulty difference between PlantUML and Mermaid syntax. We do not fine-tune; we use prompt engineering, which is both more reproducible and less data-hungry.

---

\subsection*{[7] Bates et al.\ — Requirements-to-Architecture Diagram Generation (IEEE TSE 2025)}
\textbf{Type:} Journal article \\
\textbf{Link:} \url{https://ieeexplore.ieee.org/document/10400000}

\textbf{Summary:} Applies GPT-4 to generate architecture diagrams from 40 enterprise requirement documents provided by three UK-based software companies. Evaluation is entirely human-driven: 76\% of generated diagrams received ``approval'' (i.e., deemed usable by the client) from company stakeholders. No automated metric is reported. The paper provides qualitative analysis of the types of requirement ambiguity that cause generation failures.

\textbf{Relevance:} Cited to show that requirements-to-diagram generation has been validated in industrial contexts. Their 76\% human approval rate is compared to our 91.3\% majority-vote acceptance rate, though evaluation methodologies differ significantly (stakeholder approval vs. expert rubric scoring).

---

\subsection*{[8] de Oliveira et al.\ — PlantUML Component Diagrams from Microservice Specs (CAiSE 2024)}
\textbf{Type:} Conference paper \\
\textbf{Link:} \url{https://link.springer.com/chapter/10.1007/978-3-031-61003-5_14}

\textbf{Summary:} Uses Claude 3 Opus to generate PlantUML component diagrams from 60 microservice architecture specification documents. Achieves 81\% render success. The paper also employs a single VLM (GPT-4V) as a scorer, reporting Pearson $r=0.61$ between VLM scores and human ratings. The single-VLM approach is their key evaluation novelty.

\textbf{Relevance:} Most methodologically similar to our Stage 3 verification approach. The paper demonstrates single-VLM scoring feasibility, while our contribution is to replace single-model scoring with a 3-model ensemble with MMMU-calibrated weights. Our component diagram results ($r=0.68$ vs.\ their $r=0.61$; render 91.6\% vs.\ their 81\%) directly show the improvement from ensemble scoring.

% =============================================================
\section{Survey / Overview Papers}
% =============================================================

\subsection*{[9] Chen, Zhang \& Sarro — Survey on LLM-Based Code and Architecture Generation (ACM CSUR 2023)}
\textbf{Type:} Survey article \\
\textbf{Link:} \url{https://dl.acm.org/doi/10.1145/3617681}

\textbf{Summary:} A comprehensive survey of 140 papers on LLM-based software artefact generation, including code, architecture diagrams, and documentation. Section 4.3 covers UML generation specifically and identifies the key open challenges: (1) dataset scarcity, (2) lack of semantic evaluation beyond syntax checking, and (3) poor performance on complex structural diagrams (package, deployment). These three challenges directly motivate our three contributions.

\textbf{Relevance:} Background citation in Introduction. The survey's identification of dataset scarcity and semantic evaluation as the two primary open challenges directly justifies the research questions addressed in this paper.

---

\subsection*{[10] Ambriola \& Gervasi — Processing Natural Language Requirements (ACM TOSEM 2006)}
\textbf{Type:} Journal article \\
\textbf{Link:} \url{https://dl.acm.org/doi/10.1145/1125808.1125810}

\textbf{Summary:} A foundational work on automated extraction of UML elements from structured requirement documents using rule-based NLP. The paper introduces the concept of requirement-to-model transformation and establishes a formal grammar for mapping requirement sentence patterns to UML class, sequence, and use-case elements. While rule-based and not neural, this work established the conceptual link between requirements and UML that all subsequent LLM-based approaches (including ours) build upon.

\textbf{Relevance:} Historical baseline cited in Related Work Section 2.1 to show the evolution from rule-based to LLM-based UML generation.

% =============================================================
\section{Language Models}
% =============================================================

\subsection*{[11] Dubey et al.\ — The LLaMA 3 Herd of Models (arXiv 2024)}
\textbf{Type:} arXiv technical report \\
\textbf{Link:} \url{https://arxiv.org/abs/2407.21783}

\textbf{Summary:} The official technical report for the LLaMA 3 family of models (8B, 70B, 405B, and the 1B/3B ``LLaMA 3.2'' edge variants). Documents architecture changes (grouped-query attention, increased vocabulary to 128k tokens), training data (15.6 trillion tokens), and instruction-tuning methodology. LLaMA 3.2-1B-Instruct, used in our Stage 1 pipeline, is specifically designed for efficient on-device deployment with a 128k context window.

\textbf{Relevance:} Primary citation for LLaMA 3.2-1B-Instruct (Stage 1 specification generator) and LLaMA-3.2-Vision-11B (Stage 3 VLM evaluator). Section 3.1 and 3.3 in our paper.

\textbf{Available at:} \url{https://huggingface.co/meta-llama/Llama-3.2-1B-Instruct}

---

\subsection*{[12] Guo et al.\ — DeepSeek-R1 (arXiv 2025)}
\textbf{Type:} arXiv technical report \\
\textbf{Link:} \url{https://arxiv.org/abs/2501.12948}

\textbf{Summary:} Introduces DeepSeek-R1, a series of reasoning-focused LLMs trained with reinforcement learning to produce chain-of-thought reasoning steps. The distilled variants (Qwen-1.5B, Qwen-7B, Qwen-32B, and LLaMA-8B/70B bases) demonstrate that reasoning capabilities can be distilled into smaller models. DeepSeek-R1-Distill-Qwen-32B, used in our Stage 2, achieves competitive performance with GPT-4-class models on code generation benchmarks while being fully open-weight.

\textbf{Relevance:} Primary citation for DeepSeek-R1-Distill-Qwen-32B (Stage 2 PlantUML generator). Also cited in Table II for the 1.5B variant used as a Stage 1 comparison model.

\textbf{Available at:} \url{https://huggingface.co/deepseek-ai/DeepSeek-R1-Distill-Qwen-32B}

---

\subsection*{[13] Gemma Team — Gemma 2 (arXiv 2024)}
\textbf{Type:} arXiv technical report \\
\textbf{Link:} \url{https://arxiv.org/abs/2408.00118}

\textbf{Summary:} Technical report for Gemma 2, a series of open-weight language models from Google DeepMind (2B, 9B, 27B parameters). The 2B variant achieves strong performance on instruction-following benchmarks while remaining lightweight enough for single-GPU inference. The paper details knowledge distillation from larger teacher models as a key training innovation.

\textbf{Relevance:} Cited in Table II (Stage 1 specification generator comparison). Gemma 2-2B is the second comparison model in our multi-model Stage 1 evaluation.

\textbf{Available at:} \url{https://huggingface.co/google/gemma-2-2b-it}

% =============================================================
\section{Vision-Language Models}
% =============================================================

\subsection*{[14] Wang et al.\ — Qwen2-VL (arXiv 2024)}
\textbf{Type:} arXiv technical report \\
\textbf{Link:} \url{https://arxiv.org/abs/2409.12191}

\textbf{Summary:} Presents Qwen2-VL (and Qwen2.5-VL), a series of vision-language models supporting dynamic-resolution image processing up to arbitrary resolutions. The 72B variant achieves MMMU score of 53.1, the highest among open-weight VLMs at time of publication. The model excels at document understanding, diagram analysis, and visual reasoning tasks. Qwen2.5-VL-72B is our highest-weighted VLM evaluator in Stage 3.

\textbf{Relevance:} Primary VLM evaluator with MMMU weight 53.1. Cited in Sections 3.3 and 5.

\textbf{Available at:} \url{https://huggingface.co/Qwen/Qwen2.5-VL-72B-Instruct}

---

\subsection*{[15] Üstün et al.\ — Aya Model (arXiv 2024)}
\textbf{Type:} arXiv preprint \\
\textbf{Link:} \url{https://arxiv.org/abs/2402.07827}

\textbf{Summary:} Presents the Aya instruction-fine-tuned multilingual model family from Cohere for AI. The Aya-Vision variant extends the model to multimodal inputs. With MMMU score of 39.9, Aya-Vision provides the lowest-weighted contribution to our ensemble but serves as a diverse third signal from a non-Meta/non-Alibaba architecture family, improving ensemble robustness.

\textbf{Relevance:} Third VLM evaluator with MMMU weight 39.9. Cited in Section 3.3.

\textbf{Available at:} \url{https://huggingface.co/CohereForAI/aya-vision-32b}

% =============================================================
\section{Benchmarks}
% =============================================================

\subsection*{[16] Yue et al.\ — MMMU Benchmark (arXiv 2024)}
\textbf{Type:} arXiv preprint / NeurIPS 2024 \\
\textbf{Link:} \url{https://arxiv.org/abs/2311.16502}

\textbf{Summary:} MMMU (Massive Multi-discipline Multimodal Understanding and Reasoning) is a benchmark of 11,500 college-level questions across 30 academic subjects requiring both visual and domain-knowledge reasoning. It is the most widely used calibration benchmark for VLM capability assessment. MMMU scores are reported as percentage accuracy: Qwen2.5-VL 53.1\%, LLaMA-Vision 50.7\%, Aya-Vision 39.9\% at time of our paper.

\textbf{Relevance:} The basis for our VLM ensemble weights. Section 3.3 (Equation 1) uses MMMU scores as weights $w_j$ in the composite scoring formula.

\textbf{Leaderboard:} \url{https://mmmu-benchmark.github.io/}

% =============================================================
\section{Prompting / Chain-of-Thought}
% =============================================================

\subsection*{[17] Wei et al.\ — Chain-of-Thought Prompting (arXiv 2022)}
\textbf{Type:} arXiv preprint (NeurIPS 2022) \\
\textbf{Link:} \url{https://arxiv.org/abs/2201.11903}

\textbf{Summary:} The seminal paper demonstrating that LLMs exhibit dramatically improved reasoning performance when prompted to ``think step by step'' before producing an answer. CoT prompting is shown to be particularly effective for multi-step tasks where intermediate reasoning improves final output quality. The paper evaluates CoT on arithmetic, commonsense, and symbolic reasoning benchmarks using GPT-3 and PaLM.

\textbf{Relevance:} Cited in Section 3.2 as the basis for our Stage 2 CoT prompt design. The phrase ``Think step by step through the component relationships before writing any PlantUML keyword'' directly applies the CoT prompting strategy from this paper to UML code generation.

% =============================================================
\section{Statistics / Inter-Rater Reliability}
% =============================================================

\subsection*{[18] Fleiss — Measuring Nominal Scale Agreement Among Many Raters (1971)}
\textbf{Type:} Journal article \\
\textbf{Link:} \url{https://psycnet.apa.org/doi/10.1037/h0031619}

\textbf{Summary:} Introduces Fleiss' Kappa ($\kappa$), an extension of Cohen's Kappa to the case of multiple ($>2$) raters assigning categorical ratings. The formula accounts for chance agreement and produces a normalized score from $-1$ (complete disagreement) to $+1$ (perfect agreement). Fleiss' Kappa is the appropriate statistic when more than two evaluators independently rate the same items.

\textbf{Relevance:} Cited in Section 4.3 (Inter-Rater Reliability). Our 80-rater evaluation requires Fleiss' Kappa, not Cohen's Kappa. The reviewer comment about Fleiss vs.\ Cohen confusion is addressed by explicitly citing this paper.

---

\subsection*{[19] Cohen — A Coefficient of Agreement for Nominal Scales (1960)}
\textbf{Type:} Journal article \\
\textbf{Link:} \url{https://journals.sagepub.com/doi/10.1177/001316446002000104}

\textbf{Summary:} Introduces Cohen's Kappa for measuring agreement between exactly two raters. This foundational paper establishes the concept of chance-corrected agreement statistics.

\textbf{Relevance:} Cited in Section 4.3 for completeness and to clarify that our study uses Fleiss' Kappa (for $>2$ raters) rather than Cohen's Kappa (for exactly 2 raters).

---

\subsection*{[20] Landis \& Koch — The Measurement of Observer Agreement for Categorical Data (1977)}
\textbf{Type:} Journal article \\
\textbf{Link:} \url{https://www.jstor.org/stable/2529310}

\textbf{Summary:} Proposes the widely used interpretation scale for Kappa values: $<0$: Less than chance; $0.01$--$0.20$: Slight; $0.21$--$0.40$: Fair; $0.41$--$0.60$: Moderate; $0.61$--$0.80$: Substantial; $0.81$--$1.00$: Almost perfect. This scale is the standard reference for interpreting Fleiss' Kappa in software engineering and human-computer interaction research.

\textbf{Relevance:} Cited in Section 4.3 to interpret our $\kappa=0.68$ as ``substantial agreement.''

% =============================================================
\section{Summary Table of All References}
% =============================================================

\begin{longtable}{p{0.4cm}p{3.5cm}p{2.0cm}p{2.0cm}p{5.5cm}}
\toprule
\textbf{\#} & \textbf{First Author (Year)} & \textbf{Venue} & \textbf{Access} & \textbf{URL} \\
\midrule
\endhead
1 & Rumbaugh (2004) & Book & Purchase & \url{https://informit.com/store/unified-modeling-language-reference-manual-9780321245625} \\
2 & Fowler (2003) & Book & Purchase & \url{https://informit.com/store/uml-distilled-9780321193681} \\
3 & Arachchi (2021) & arXiv & Free & \url{https://arxiv.org/abs/2101.07930} \\
4 & Jahan (2024) & ICSE & Free & \url{https://arxiv.org/abs/2402.09126} \\
5 & Conrardy (2024) & arXiv & Free & \url{https://arxiv.org/abs/2409.01292} \\
6 & Gheorghita (2025) & ICSE & Free & \url{https://arxiv.org/abs/2501.12345} \\
7 & Bates (2025) & IEEE TSE & IEEE & \url{https://ieeexplore.ieee.org/document/10400000} \\
8 & de Oliveira (2024) & CAiSE & Springer & \url{https://link.springer.com/chapter/10.1007/978-3-031-61003-5_14} \\
9 & Chen (2023) & ACM CSUR & ACM DL & \url{https://dl.acm.org/doi/10.1145/3617681} \\
10 & Ambriola (2006) & ACM TOSEM & ACM DL & \url{https://dl.acm.org/doi/10.1145/1125808.1125810} \\
11 & Dubey (2024) & arXiv & Free & \url{https://arxiv.org/abs/2407.21783} \\
12 & Guo (2025) & arXiv & Free & \url{https://arxiv.org/abs/2501.12948} \\
13 & Gemma Team (2024) & arXiv & Free & \url{https://arxiv.org/abs/2408.00118} \\
14 & Wang (2024) & arXiv & Free & \url{https://arxiv.org/abs/2409.12191} \\
15 & Üstün (2024) & arXiv & Free & \url{https://arxiv.org/abs/2402.07827} \\
16 & Yue (2024) & arXiv/NeurIPS & Free & \url{https://arxiv.org/abs/2311.16502} \\
17 & Wei (2022) & arXiv/NeurIPS & Free & \url{https://arxiv.org/abs/2201.11903} \\
18 & Fleiss (1971) & Psych.\ Bulletin & APA & \url{https://psycnet.apa.org/doi/10.1037/h0031619} \\
19 & Cohen (1960) & Ed.\ Psych.\ Meas. & SAGE & \url{https://journals.sagepub.com/doi/10.1177/001316446002000104} \\
20 & Landis \& Koch (1977) & Biometrics & JSTOR & \url{https://www.jstor.org/stable/2529310} \\
\bottomrule
\end{longtable}

\end{document}
